\documentclass[11pt]{article}
\pdfinfoomitdate=1\pdftrailerid{}
\usepackage[letterpaper,margin=.7in,headheight=15pt,headsep=18pt,footskip=24pt]{geometry}
\usepackage[T1]{fontenc}
\usepackage{lmodern,amsmath,amssymb,array,booktabs,fancyhdr,tikz}
\usetikzlibrary{arrows.meta}
\renewcommand{\familydefault}{\sfdefault}
\setlength{\parindent}{0pt}\setlength{\parskip}{6pt}
\pagestyle{fancy}\fancyhf{}
\fancyhead[L]{\small Week 66 / Lamplighter streets / K--1 adult guide}
\fancyfoot[L]{\scriptsize Bellingham Math Circle / Week 66 / GGT66-K-FAC-v2}
\fancyfoot[R]{\thepage}
\renewcommand{\headrulewidth}{0pt}
\newcommand{\head}[1]{\par\textbf{\large #1}\par}
\newcommand{\sub}[1]{\par\textbf{#1}\par}
\newcommand{\keyrow}[8]{#2 & #3 & #4 & \texttt{#5} & $#6+#7$ & \textbf{#8}\\}
\newcommand{\allways}[3]{\textbf{#2 ways:} \texttt{#3}}
\newcommand{\exactword}[2]{\texttt{#2}}
\newenvironment{keytable}{\par\smallskip\begingroup\small\setlength{\tabcolsep}{7pt}\renewcommand{\arraystretch}{1.15}\begin{tabular}{@{}llllcr@{}}\textbf{Picture}&\textbf{Lamps on}&\textbf{Walker}&\textbf{One shortest word}&\textbf{Steps+flips}&\textbf{Fewest}\\\midrule}{\end{tabular}\endgroup\par\smallskip}
\newcommand{\demo}[3]{%
\begin{tikzpicture}[x=4.6mm,y=4.6mm,baseline=0pt]
\draw[gray] (-3,0)--(3,0);
\foreach \p in {-3,-2,...,3}{\draw[fill=white] (\p,0) circle (.23);}
\ifnum#1=1\fill (3,0) circle (.23);\fi
\begin{scope}[shift={(#2,.34)},scale=.11]
\draw[fill=black!45] (-2.5,0)--(2.5,0)--(2.5,.9)--(1.8,.9)
.. controls (1.2,2.2) and (.9,3.2) .. (.9,4)--(-.9,4)
.. controls (-.9,3.2) and (-1.2,2.2) .. (-1.8,.9)--(-2.5,.9)--cycle;
\draw[fill=black!45] (0,5.35) circle (1.45);
\end{scope}
\draw (0,-.38)--(.22,-.56)--(.22,-.82)--(-.22,-.82)--(-.22,-.56)--cycle;
\ifnum#3>0
\foreach \i in {1,...,#3}{\draw[fill=black!15] (4.2,{.48*(\i-1)}) rectangle +(0.48,.48);}
\fi
\end{tikzpicture}}
\newcommand{\record}[1]{\begin{tikzpicture}[x=3mm,y=3mm,baseline=0pt]
\foreach \i in {1,...,4}{\ifnum\i>#1\def\shade{white}\else\def\shade{black!35}\fi
\draw[fill=\shade] (\i,0) rectangle +(1,1);}\end{tikzpicture}}
\begin{document}
\head{Mathematical overview: a picture has a shortest build}
\textbf{Fewest moves = required flips + shortest visiting walk.} Start all-off with the walker home. A step and a flip under the walker are separate moves. Every lamp ending on needs a flip and a visit; the walker must also reach its pictured finish. Conversely, a shortest visiting walk, flipping each required lamp once, meets this lower bound. Point to the unavoidable flips and trace the walk: a short trial alone does not prove optimality.

\textbf{On a line, compare the two visiting orders.} Reach the farthest lit lamp on each side of home (use home if that side has none), one end then the other, and then the finish. The cheaper order gives the shortest walk. Ending at an end lamp favors visiting the other end first (Problem 3); moving the finish can tie or reverse the best order (Problems 2 and 4).

\textbf{Different shortest ways.} Flip each required lamp exactly once along a shortest walk. When that walk stands on a lamp twice, either visit may carry its flip: Problem 5 has 2 and 4 shortest ways.

\textbf{Exact counts and neighboring pictures.} If the minimum is $d$, exactly $d,d+2,d+4,\ldots$ moves work. A flip and flip back adds two. Conversely, lit lamps plus the walker's steps from home changes evenness at every move, so all routes to one picture have the same evenness. Adding two explains possibility, not exclusion of other counts. One-move neighbors differ in minimum by at most one (append that move in either direction), and parity excludes a tie: they differ by \emph{exactly one}. Thus every right-hand count in Problems 6--7 is one above or below the left count. Problem 7's $9\to8,8,8$ is a dead end \emph{for distance}, not for movement; Problem 6 contrasts with it.

\textbf{Scope and destination.} These facts hold on a line with finitely many lit lamps; the seven-lamp street contains all printed shortest routes. In the first hour K--1 finds that finishing place and visiting order change the count. Later routes compare complete lists, a dead end and evenness. A build is an experiment, a proposed minimum a conjecture, and forced flips/walking an explanation. General dead-end classification belongs to the older tables: see ``Mathematical destination,'' base guide page 1.

\head{Preparation and adult support}
\textbf{Status.} Unpiloted and physically unrehearsed: counter/pawn fit, flipping, cube stacks, pencil records, staffing and timing remain untested.

\textbf{Per pair.} Use the week's kit: nine two-sided lamp counters no wider than 17 mm (or paper squares dark on one side), one visibly different pawn, about 24 cubes, and pencils with erasers. Set out 7 counters; keep two spares. Use linking cubes if available. The K--1 lamps are 20 mm circles on 24 mm centers, with separate 17 mm pawn spots above; choose a pawn whose base fits a spot. Dark side up means on, matching filled circles. Reset by saying ``all lamps off'' (every counter light side up); avoid ``dark street.'' Counters stay on their lamps; the pawn never covers a counter.

\textbf{Print and prepare.} Print eight pages per pair from the current draft \texttt{k-1.pdf}: US Letter, single-sided, Actual Size. After approval, use \texttt{week-66-k-1.pdf} in the Week 66 print folder. Each page has a working street; small pictures are targets. Older tables keep \texttt{week-66-students.pdf} and the base guide. Supply real counters, pawn and cubes; no materials PDF prints them. Allow 5--10 extra minutes for printing and two pair kits; ready-made counters need no cutting.

\textbf{Prerequisites and staffing.} Children match two-part pictures and count/compare towers to about 10; longer tries are possible. No independent reading, signed numbers or letter writing is needed. One adult anchors two K--1 pairs, reads and records on request. The mover chooses steps; the flipper points to the pawn's lamp before flipping. Whoever acts adds a cube; swap jobs between pictures. Point-before-flip checks legality; the adult spot-checks and watches resets. If adults must direct every action, stay with one picture.
\newpage
\head{Launch, first route and the first two problems}
\textbf{Join the whole-group launch.} Keep the base guide's short demonstration of separate steps and flips, and its comparison of two routes to one goal. Let children handle counters first. Back at the K--1 table, practise a \emph{non-task try}: all off, pawn in the spot above the rightmost lamp. Flip there and add a cube; step one place left and add another. The lit lamp stays put. Colour two squares on a scrap strip for the two cubes, then dismantle the tower and reset all-off, pawn home. Problems 1--4 record a best try this way in the squares under each picture; leave dot/arrow records until Problem 5.
\begin{center}
\small
\begin{tabular}{c@{\quad}c@{\quad}c}
Input: no cubes & After a flip: 1 cube & After a step left: 2 cubes\\[3pt]
\demo{0}{3}{0}&\demo{1}{3}{1}&\demo{1}{2}{2}\\[6pt]
\record{0}&\record{1}&\record{2}
\end{tabular}
\end{center}
Before \emph{every target attempt}, reset all-off, pawn home, tower apart; keep the target visible. Record a new best before reusing cubes. A single pencil stroke per square is enough; let the adult erase surplus marks when a shorter try wins. Twelve squares are capacity, not a promised answer or a limit: note longer tries beside them.

\textbf{A one-hour first visit.} Allow about 5 minutes for the common launch, 3 for table practice, 10--12 each for Problems 1 and 2, a 3-minute stretch, then about 15 for selected Problem 3 or 4 pictures and 5 for sharing an unavoidable-moves argument. Leave transition time. Stop earlier if matching and replaying are still absorbing. Save Problem 5's complete lists and Problems 6--8 for a later visit or a ready pair with adult attention. For that pair, take 6 and 7 together; 8 can be a separate return investigation. Finishing every page is not the goal.

\textbf{How to read the key.} Each problem is on the same-numbered student page. TL/TR/BL/BR mean top-left/top-right/bottom-left/bottom-right. Adult shorthand only: H = home; 2L = second lamp left of home; 3R = third right. Words are read left to right: L/R is one step, F is one flip. Children point and act; Problem 5 uses dot/arrow records. The table gives one shortest route unless ``every way'' is requested.

\sub{Problem 1 (student page 1): separate walking from flipping}
\begin{keytable}
\keyrow{1TL}{TL}{H}{H}{F}{0}{1}{1}
\keyrow{1TR}{TR}{none}{2L}{LL}{2}{0}{2}
\keyrow{1BL}{BL}{2R}{2R}{RRF}{2}{1}{3}
\keyrow{1BR}{BR}{3L}{3L}{LLLF}{3}{1}{4}
\end{keytable}
\textbf{Reasoning.} The home lamp needs only a flip. The all-off picture needs only two steps. The other two force a walk to the indicated lamp and one flip there. Each displayed word meets these lower bounds.
\textbf{Hints, only if needed:} ``Which lamp needs a flip?'' ``Does the walker match the picture too?'' A child who flips a lamp in the all-off picture may undo it legally, then compare the longer tower.

\sub{Problem 2 (student page 2): same lamp, different finish}
\begin{keytable}
\keyrow{2TL}{TL}{2R}{3R}{RRFR}{3}{1}{4}
\keyrow{2TR}{TR}{2R}{1R}{RRFL}{3}{1}{4}
\keyrow{2BL}{BL}{2R}{H}{RRFLL}{4}{1}{5}
\keyrow{2BR}{BR}{2R}{1L}{RRFLLL}{5}{1}{6}
\end{keytable}
\textbf{Reasoning.} All four need one flip at 2R. Reaching it takes two steps; the finish then needs one, one, two or three more. The top pair tie despite ending in different places.
\textbf{Hint:} ``What changes between these pictures, and what stays the same?'' \textbf{Continuation:} ask which new ending would make the same lit-lamp picture easiest; ending on that lamp needs just the outward walk and its flip.
\newpage
\head{Visiting order and finding every shortest way}
\sub{Problem 3 (student page 3): which end first?}
\begin{keytable}
\keyrow{3TL}{TL}{1R,2R}{2R}{RFRF}{2}{2}{4}
\keyrow{3TR}{TR}{1L,2R}{2R}{LFRRRF}{4}{2}{6}
\keyrow{3BL}{BL}{3L,1R}{3L}{RFLLLLF}{5}{2}{7}
\keyrow{3BR}{BR}{2L,3R}{3R}{LLFRRRRRF}{7}{2}{9}
\end{keytable}
\textbf{Reasoning.} Two lamps force two flips. TL needs two right steps. For the other pictures, the walk must reach both lit ends and finish at the pictured end. Going first to the \emph{other} end avoids walking out to the finishing end and back again: TR has 4 steps rather than 8; BL has 5 rather than 11; BR has 7 rather than 13. These are walking counts before the two flips.
\textbf{Hints:} ``Where must the walker finish?'' ``Can you trace both orders without flipping yet?'' Do not announce the preferred order before children have a route to compare.

\sub{Problem 4 (student page 4): changing the finish can change the best order}
\begin{keytable}
\keyrow{4TL}{TL}{2L,1R}{1R}{LLFRRRF}{5}{2}{7}
\keyrow{4TR}{TR}{2L,1R}{H}{LLFRRRFL}{6}{2}{8}
\keyrow{4BL}{BL}{2L,1R}{1L}{RFLLLFR}{5}{2}{7}
\keyrow{4BR}{BR}{2L,1R}{2L}{RFLLLF}{4}{2}{6}
\end{keytable}
\textbf{Reasoning.} Always two flips. Trace the trip to 2L first, then 1R and the finish; compare going to 1R first. The walking counts (left first/right first) are TL 5/7, TR 6/6, BL 7/5, BR 8/4. The cheaper route gives the table. At home either order ties; this target has exactly two shortest words, \texttt{LLFRRRFL} and \texttt{RFLLLFRR}.
\textbf{Hint:} ``Which lit lamp will you go to first? What if you went to the other one first?'' \textbf{Continuation:} compare the two equal home-ending routes by replay, not just their tower heights.

\sub{Problem 5 (student page 5): every way, not merely one}
Before starting, replay the page's small dot-then-right-arrow example: all off at home, flip home, then step right. Read input, intermediate and finish. Its two moves begin one of the top picture's answers. That is fine for ``every way,'' but replay it to the side and do not continue it. Children mark moves in order, one per square, and leave unused end squares blank. The adult can draw requested marks, but the child should replay and compare rows. A fifth blank row does not promise five solutions.

\textbf{Top picture: home and 1R lit, walker home.} Fewest \textbf{4}: two steps and two flips. \allways{5top}{2}{FRFL, RFLF}. The walk must go right and return. The 1R lamp can only be flipped at the turn; home can be flipped before leaving or after returning. These are all choices.

\textbf{Bottom picture: home, 1R and 2R lit, walker home.} Fewest \textbf{7}: four steps and three flips. \allways{5bottom}{4}{FRFRFLL, FRRFLFL, RFRFLLF, RRFLFLF}. The walk must be RRLL. Flip 2R at the turn; choose whether to flip 1R outward or homeward, and home before leaving or after returning. Two independent choices give all four ways. A route with fewer steps cannot visit 2R and return; fewer flips cannot light all three.

\textbf{Hints:} ``Which lamps do you stand on more than once?'' ``When does each of your ways light the home lamp: before the walker leaves, or after it comes back?'' Later: ``Could you leave this lamp off now and light it when you return?'' Readiness means the child can recognize a replayed route; if only the adult can compare the records, return to improving one route.
\newpage
\head{One more move, and exact-size towers}
\sub{Problem 6 (student page 6): a harder next picture exists}
The arrows describe how pictures differ; they are \emph{not} the route to use. Solve all four pictures afresh from all-off at home. The right-hand pictures are top = L, middle = R, bottom = F; the diagonal arrow labeled L is still a \emph{left step}, not a diagonal move on the street.
\begin{keytable}
\keyrow{6start}{Left}{2L,1L}{1L}{LFLFR}{3}{2}{5}
\keyrow{6L}{Top right (L)}{2L,1L}{2L}{LFLF}{2}{2}{4}
\keyrow{6R}{Middle right (R)}{2L,1L}{H}{LFLFRR}{4}{2}{6}
\keyrow{6F}{Bottom right (F)}{2L}{1L}{LLFR}{3}{1}{4}
\end{keytable}
\textbf{Answer and reason.} The middle-right R picture needs the most: \textbf{6}. Each route must reach 2L; the ending forces respectively 3, 2, 4 and 3 steps, plus the shown flips. Stepping toward home from the left picture creates a target needing more walking from the original start.
\textbf{Hint:} ``Where did this attempt start? Where must it finish?''

\sub{Problem 7 (student page 7): every next picture is easier}
\begin{keytable}
\keyrow{7start}{Left}{1L,H,2R}{H}{FLFRRRFLL}{6}{3}{9}
\keyrow{7L}{Top right (L)}{1L,H,2R}{1L}{FRRFLLLF}{5}{3}{8}
\keyrow{7R}{Middle right (R)}{1L,H,2R}{1R}{FLFRRRFL}{5}{3}{8}
\keyrow{7F}{Bottom right (F)}{1L,2R}{H}{LFRRRFLL}{6}{2}{8}
\end{keytable}
\textbf{Answer and reason.} The \emph{left} picture needs the most: \textbf{9}. Visiting 1L and 2R and returning home needs six steps plus three flips. Finishing one step either side saves one walking step; turning home off saves one flip but no walking. All three alternatives need \textbf{8}, so all possible next moves make the target easier. Appending a move to a nine-move route gives a ten-move attempt, not the new minimum.
\textbf{Hints:} ``Are we including the left picture?'' (The three right pictures alone tie at 8.) ``What must happen if we finish at home?'' The four towers would use 33 cubes; the shaded squares retain the counts while the pair reuses its 24 cubes. A ready pair can invent a similar target and compare all three neighbors; the base guide supplies the general rule.

\sub{Problem 8 (student page 8): which exact counts work?}
\textbf{Top row: circle 1 and 3; not 2 or 4.} Examples: \exactword{8top}{F}, \exactword{8top}{FFF}.\par
\textbf{Middle row: circle 2 and 4; not 1 or 3.} Examples: \exactword{8middle}{LR}, \exactword{8middle}{LRLR}. The picture is already built in zero moves, but no zero tower is offered.\par
\textbf{Bottom row: circle 4 and 6; not 5 or 7.} Examples: \exactword{8bottom}{RFLL}, \exactword{8bottom}{RFLLFF}.

\textbf{Reasoning and hints.} All successful exact counts are the minimum plus an even number (overview). Let the child experiment, then ask ``Can two extra moves leave the same picture?'' and ``What do the working tower heights have in common?'' Adding a cancelling pair explains the possible counts. For the impossible counts, the adult may pair up lit-lamp counters together with one cube per step the pawn is from home: every move switches between an even and an odd total. A few failed attempts are not proof. Keep the parity argument optional, not a condition for enjoying the investigation.

\textbf{Where it goes.} Once the K--1 pages are approved, this section replaces only ``K--1 tables.'' in ``Materials and preparation,'' page 1 of \texttt{week-66-facilitator.pdf} (GGT66-FAC-v2); retain its shared launch and older-band guide. In its page-1 ``What children discover'' paragraph, add: ``K--1 uses the separate picture packet and this section to compare finishing places and visiting orders before optional route lists and dead ends.'' The approved student file will be \texttt{lowell-math-circle-year-2/}\allowbreak\texttt{week-66/week-66-k-1.pdf}.
\end{document}
